\section{Fundamentos da Movimentação Dentária sob a Ótica Computacional}

A movimentação dentária ortodôntica (MDO) é, em sua essência, um fenômeno mecano-transdutivo de alta complexidade. Trata-se de um processo de remodelação tecidual orquestrado que envolve a adaptação do ligamento periodontal (LPD), do cemento radicular e do osso alveolar frente a estímulos mecânicos exógenos. Quando um vetor de força é aplicado a uma coroa dentária, as fibras colágenas do LPD experimentam gradientes de tensão (tração) e de pressão (compressão). Este gradiente de deformação (\textit{strain}) altera o fluxo de fluidos intersticiais e deforma o citoesqueleto de fibroblastos, osteoblastos e osteócitos, deflagrando uma cascata catabólica e anabólica mediada por citocinas, prostaglandinas e fatores de crescimento (como RANKL e OPG)~\cite{zhang2024recursive}.

Do ponto de vista da física do estado sólido e da biomecânica computacional, o complexo dente-LPD-osso alveolar não se comporta de maneira linear. Pelo contrário, \destaque{o periodonto atua como um sistema viscoelástico, anisotrópico e hiperelástico}, cuja resposta fenomenológica é estritamente dependente da magnitude (forças leves contínuas versus forças pesadas intermitentes), direção, duração e frequência da carga aplicada~\cite{khoshfekr2025digital}. Para que um gêmeo digital seja epistemologicamente válido, esse comportamento biomecânico não-linear deve ser capturado mediante o uso de equações constitutivas sofisticadas que modelem a histerese tecidual e o relaxamento de tensões no domínio do tempo. Tais modelos matemáticos permitem que os algoritmos de \textit{solver} antecipem, com considerável acurácia tridimensional, não apenas o deslocamento inicial imediato (fase elástica), mas também a complexa fase de retardo e a fase de movimentação progressiva que ocorrem nas semanas subsequentes à ativação ortodôntica.

A Ortodontia apresenta uma topologia de dados que a torna um terreno excepcionalmente fértil para a modelagem preditiva \textit{in silico}. Em primeiro lugar, a morfologia estática do sistema -- arquitetura radicular, espessura cortico-esponjosa e limites anatômicos maxilomandibulares -- pode ser digitalizada com resolução submilimétrica através da Tomografia Computadorizada de Feixe Cônico (CBCT) fundida a varreduras de superfície intraoral (IOS)~\cite{katsoulakis2024digital}. Em segundo lugar, e de suma importância para as condições de contorno de um modelo matemático, os vetores e os momentos de força aplicados pelos aparelhos (fios de NiTi, elásticos, alinhadores termoplásticos) são quantificáveis, controláveis e previsíveis. Por fim, o lapso temporal inerente à biologia óssea (onde a remodelação clinicamente observável leva de semanas a meses) confere ao sistema computacional uma janela de oportunidade ímpar para o processamento de dados robustos, calibração algorítmica e validação prospectiva iterativa dos modelos matemáticos construídos~\cite{duggal2026exploring}.
